% GCMM 2026 poster (Paper 59), layout of the ICAISC 2026 poster — adapted from W. Bożejko's PWr poster template
% (ICCS 2023, "Determination of the lower bounds...", Warsaw beamer
% theme, green structure colour, Times font, 3-column block layout).

\documentclass[final]{beamer}
  \mode<presentation>
  {
      \usetheme{Warsaw}
  }

% --- Cyberpunk neon palette ----------------------------------------
\definecolor{cyberviolet}{HTML}{6A1B9A}   % deep violet  (block headers)
\definecolor{cybermagenta}{HTML}{F72585}  % neon magenta (Fibonacci / alert)
\definecolor{cybercyan}{HTML}{4CC9F0}     % neon cyan    (accents)
\usecolortheme[named=cyberviolet]{structure}

% block headers: violet bar, white text, faint violet body tint
\setbeamercolor{block title}{fg=white,bg=cyberviolet}
\setbeamercolor{block body}{fg=black,bg=cyberviolet!6}
% alert (key-result) block: neon magenta
\setbeamercolor{block title alerted}{fg=white,bg=cybermagenta}
\setbeamercolor{block body alerted}{fg=black,bg=cybermagenta!10}
% header band + title text
\setbeamercolor{headline}{fg=white,bg=cyberviolet}
\setbeamercolor{title in headline}{fg=white}
\setbeamercolor{author in headline}{fg=cybercyan}
\setbeamercolor{institute in headline}{fg=white}

\usepackage{times}
\usepackage{amsmath,amssymb}
\usepackage[english]{babel}
\usepackage[orientation=portrait,size=a0,scale=1.0]{beamerposter}
\usepackage[utf8]{inputenc}
\usepackage[T1]{fontenc}
\usepackage{graphicx}
\usepackage{booktabs}
\usepackage{tikz}
\usepackage{qrcode}
\definecolor{cadj}{HTML}{5B7DB1}
\definecolor{cfib}{HTML}{9A342D}
\definecolor{cdyn}{HTML}{E08A1E}
\definecolor{cmot}{HTML}{2A9D8F}
\newlength{\colh}\setlength{\colh}{0.83\textheight}
\newcommand{\fig}[2][0.98]{\begin{center}\includegraphics[width=#1\columnwidth]{img/#2}\end{center}}
\usetikzlibrary{positioning,arrows.meta}

\title{Gate-Model Quantum Neighborhood Search for the Permutation Flow Shop Problem via Fixed-Angle QAOA}

\author[Wojew\'odzki \& Bo\.{z}ejko]
  {Remigiusz Wojew\'odzki$^1$ \and Wojciech Bo\.{z}ejko$^1$}

\institute[Wroc\l aw University of Science and Technology]
  {$^1$\,Department of Control and Quantum Computing,
    Wroc\l aw University of Science and Technology, Poland\\[1ex]
    \texttt{remigiusz.wojewodzki@pwr.edu.pl} \quad
    \texttt{wojciech.bozejko@pwr.edu.pl}}

\date{GCMM 2026, Wroc\l aw, Poland}

\begin{document}
\begin{frame}
\frametitle{
  \begin{columns}[T]
    \begin{column}{.01\paperwidth}\end{column}
    \begin{column}{.10\paperwidth}
      \begin{center}
        \setlength{\fboxsep}{6mm}\colorbox{white}{\includegraphics[height=11cm]{logo_PWR_square.png}}
      \end{center}
    \end{column}
    \begin{column}{.86\paperwidth}
      \begin{center}
        \usebeamercolor{title in headline}{\color{fg}\textbf{\veryHuge{\inserttitle}}\\[1ex]}
        \usebeamercolor{author in headline}{\color{fg}\LARGE{\insertauthor}\\[1ex]}
        \usebeamercolor{institute in headline}{\color{fg}\large{\insertinstitute}\\[1ex]}
      \end{center}
    \end{column}
    \begin{column}{.01\paperwidth}\end{column}
  \end{columns}
}

{\large
\begin{columns}[T]
\setlength{\columnsep}{45mm}

% =====================================================================
% LEFT COLUMN: the method
% =====================================================================
\column{0.31\textwidth}
\begin{minipage}[t][\colh][s]{\linewidth}

\begin{block}{One QUBO, two quantum back-ends}
\vspace*{2mm}
\textbf{PFSP:} $n$ jobs visit $m$ machines in one common order, and we look
for the order with the smallest makespan $C_{\max}$. Local search improves the
order by \textbf{moves}, i.e.\ sets of job swaps; a \textbf{neighborhood} is the
family of moves the search may try. Earlier work wrote the choice
of a move as a \textbf{QUBO} and solved it on a quantum annealer; here the
\textbf{same QUBOs} run on \textbf{gate-model} chips through QAOA.
\vspace{2mm}

{\normalsize\textbf{QUBO}: a quadratic optimisation over 0/1 variables.
\textbf{QAOA}: a short quantum circuit whose samples are good 0/1 answers.
\textbf{PRD}: percent above the Taillard lower bound,
$(C_{\max}-\mathrm{LB})/\mathrm{LB}\cdot100\%$.}
\vspace{4mm}
\begin{center}
\resizebox{0.97\columnwidth}{!}{%
\begin{tikzpicture}[
    box/.style={draw,rounded corners,align=center,minimum height=50mm,line width=2pt},
    >={Stealth[length=8mm]}, line width=2pt]
  \node[box,fill=cyberviolet!15,minimum width=80mm] (ls) at (0,0) {\textbf{ILS / SA}\\current order};
  \node[box,fill=cybermagenta!20,minimum width=110mm] (q) at (16,0) {\textbf{choice of move}\\\textbf{as QUBO}};
  \node[box,fill=black!8,minimum width=110mm,minimum height=36mm] (a) at (34,4) {quantum annealer\\earlier work};
  \node[box,fill=cyberviolet,text=white,minimum width=110mm,minimum height=36mm] (g) at (34,-4) {\textbf{gate-model QAOA}\\this work};
  \draw[->] (ls)--(q);
  \draw[->] (q.east)--(a.west);
  \draw[->] (q.east)--(g.west);
  \draw[->,dashed] (g.south)--++(0,-2.5)-|(ls.south) node[pos=.25,below]{new order};
\end{tikzpicture}}
\end{center}
\vspace*{2mm}
\end{block}

\vfill

\begin{block}{Four neighborhoods, four conflict graphs}
\vspace*{2mm}
\fig[0.7]{p_conflict_graphs.png}
Each candidate swap is one qubit and each clashing pair is one two-qubit
gate.
Fibonacci is a sparse chain; in Adjacent every pair clashes, and Dynasearch and Motzkin are dense.
\vspace*{2mm}
\end{block}

\vfill

\begin{block}{One circuit per move, angles frozen}
\vspace*{2mm}
\fig[0.85]{p_landscape1.png}
The two QAOA angles are tuned once on a simulator (star) and then frozen, so
every move on hardware is one submission. Each circuit is sampled 4096 times
and the sample with the lowest QUBO value becomes the move.
\vspace*{2mm}
\end{block}

\vfill

\begin{block}{Three chip layouts, same four QUBOs}
\vspace*{2mm}
\fig{p_coupling_maps.png}
A chip can apply a two-qubit gate only along its links, so the compiler
inserts extra gates for every missing link.
\vspace*{2mm}
\end{block}

\end{minipage}

% =====================================================================
% MIDDLE COLUMN: compilation and IBM hardware
% =====================================================================
\column{0.31\textwidth}
\begin{minipage}[t][\colh][s]{\linewidth}

\begin{block}{Circuit size of one move after compilation}
\vspace*{2mm}
\fig[1.0]{p_depth.png}
Log scale; CZ is the chip's two-qubit gate. Star~16 counts include MOVE operations
(CZ alone: Fibonacci 14, Motzkin 576, Dynasearch 696).
\vspace*{2mm}
\end{block}

\vfill

\begin{alertblock}{Circuit size: same order on every chip}
\vspace*{2mm}
\begin{center}
{\large\mbox{Fibonacci $\ll$ Adjacent $\ll$ Motzkin $<$ Dynasearch}}
\end{center}
\vspace{1mm}
A deeper circuit collects more noise. Compiled size ranks the neighborhoods
identically on all four chips, the same order in which the neighborhoods fit
on the quantum annealer of our earlier work.
\vspace*{2mm}
\end{alertblock}

\vfill

\begin{block}{IBM Heron r2: one move, ten instances}
\vspace*{2mm}
\fig[0.95]{p_hw_prd1.png}
\begin{center}
{\LARGE \textbf{27.9\,\%} $\rightarrow$ \textbf{22.1--25.0\,\%}}
\end{center}
Mean PRD before and after one move; 80 runs, about 460~s of QPU time. ILS and SA differ by less than half a
point, because one move from the same start is the same computation.
\vspace*{2mm}
\end{block}

\vfill

\begin{block}{Why only one move: the queue}
\vspace*{2mm}
\fig[0.85]{p_hw_time.png}
A move needs 3--7~s of QPU time but took 10--288~s on the clock, so no run
reached a second move within the 10~s budget.
\vspace*{2mm}
\end{block}

\end{minipage}

% =====================================================================
% RIGHT COLUMN: comparison and conclusions
% =====================================================================
\column{0.31\textwidth}
\begin{minipage}[t][\colh][s]{\linewidth}

\begin{block}{Gate model, annealer and classical search (ILS)}
\vspace*{2mm}
\fig{p_compare.png}
One gate-model move is compared with 10~s of annealer or classical search, so
the quantum back-ends are not ranked: the gate move is lowest on Dynasearch,
the full annealer on the other three. No quantum back-end comes close to
classical ILS (about 3\,\%) at $n=20$.
\vspace*{2mm}
\end{block}

\vfill

\begin{block}{Same circuits on three chips}
\vspace*{2mm}
\fig[0.95]{p_arch1.png}
On the first instance, Adjacent and Fibonacci give almost the same result on
every chip, while Dynasearch and Motzkin change from chip to chip.
\vspace*{2mm}
\end{block}

\vfill

\begin{alertblock}{Short circuits agree across chips, long ones do not}
\vspace*{2mm}
\begin{center}
\renewcommand{\arraystretch}{1.3}
\begin{tabular}{lrp{0.55\columnwidth}}
\textbf{Fibonacci} & 14 CZ   & same $C_{\max}$ on 6 of 9 instances, rest within 0.85 points\\
\textbf{Adjacent}  & 195 CZ  & same $C_{\max}$ on 9 of 9 instances\\
\textbf{Motzkin}   & 1455 CZ & PRD 16.96 on Heron~r2, 13.31 on Star~16 (one instance)\\
\end{tabular}
\end{center}
\vspace*{1mm}
CZ = two-qubit gate, counted on Heron~r2; instances are those that fit
16 qubits. With one instance for the deep circuits, the chip and sampling
noise cannot be separated yet.
\vspace*{2mm}
\end{alertblock}

\vfill

\begin{block}{Conclusions}
\vspace*{2mm}
\begin{itemize}
  \item One QUBO per neighborhood drives both a quantum annealer and a QAOA
    circuit with frozen angles.
  \item One QAOA move lowers PRD from 27.9\,\% to 22.1--25.0\,\%; classical
    ILS is still far better at $n=20$.
  \item Compiled circuit size ranks the neighborhoods the same way on every
    chip; a move is one submission and about 30~s on the clock, and the
    number of billed circuits grows linearly with $n$.
  \item Next: the full protocol (10 instances $\times$ 5 seeds $\times$
    6 budgets) on every chip, and larger $n$.
\end{itemize}
\vspace*{2mm}
\end{block}

\vfill

\begin{block}{Data and selected references}
\vspace*{1mm}
\begin{minipage}[c]{0.24\columnwidth}\qrcode[height=5cm]{https://doi.org/10.5281/zenodo.24106982}\end{minipage}\hfill
\begin{minipage}[c]{0.73\columnwidth}Code, frozen angles and all hardware runs:\\
{\small\texttt{doi.org/10.5281/zenodo.24106982}}\end{minipage}\\[1.2ex]
{\small
[1] Farhi, E., Goldstone, J., Gutmann, S. \emph{A quantum approximate
    optimization algorithm.} arXiv:1411.4028 (2014).\\[0.4ex]
[2] Bo\.zejko, W., Uchro\'nski, M., Wodecki, M. \emph{The new exponential
    Fibonacci neighborhood as QUBO\ldots} Sci.\ Rep.\ (2026).\\[0.4ex]
[3] Wojew\'odzki, R., Bo\.zejko, W. \emph{Scalable quantum neighborhood
    search for the PFSP via windowed QUBO decomposition.} SSRN (2026).
}
\vspace*{1mm}
\end{block}

\end{minipage}
\end{columns}
}
\end{frame}
\end{document}
